\documentclass{tufte-handout}

%\geometry{showframe}
%\geometry{showframe}% for debugging purposes -- displays the margins

\usepackage{amsmath}
\usepackage{amssymb}
\usepackage{cleveref}
\crefname{Proposition}{Proposition}{Propositions}
\Crefname{Proposition}{Proposition}{Propositions}
\crefname{Theorem}{Theorem}{Theorems}
\Crefname{Theorem}{Theorem}{Theorems}
\crefname{Definition}{Definition}{Definitions}
\Crefname{Definition}{Definition}{Definitions}
\crefname{Corollary}{Corollary}{Corollaries}
\Crefname{Corollary}{Corollary}{Corollaries}
\crefname{Lemma}{Lemma}{Lemmas}
\Crefname{Lemma}{Lemma}{Lemmas}
\crefname{Example}{Example}{Examples}
\Crefname{Example}{Example}{Examples}
\usepackage{amsthm}

% Set up the images/graphics package
\usepackage{graphicx}
\setkeys{Gin}{width=\linewidth,totalheight=\textheight,keepaspectratio}
\graphicspath{{graphics/}}

% The following package makes prettier tables.  We're all about the bling!
\usepackage{booktabs}

% The units package provides nice, non-stacked fractions and better spacing
% for units.
\usepackage{units}

% The fancyvrb package lets us customize the formatting of verbatim
% environments.  We use a slightly smaller font.
\usepackage{fancyvrb}
\fvset{fontsize=\normalsize}

% Small sections of multiple columns
\usepackage{multicol}

% Provides paragraphs of dummy text
\usepackage{lipsum}

% Defines colors
\usepackage{xcolor}
\definecolor{blue}{cmyk}{0.63, 0.37, 0, 0.57}     
\definecolor{ltblue}{RGB}{78,150,179}

% These commands are used to pretty-print LaTeX commands
\newcommand{\doccmd}[1]{\texttt{\textbackslash#1}}% command name -- adds backslash automatically
\newcommand{\docopt}[1]{\ensuremath{\langle}\textrm{\textit{#1}}\ensuremath{\rangle}}% optional command argument
\newcommand{\docarg}[1]{\textrm{\textit{#1}}}% (required) command argument
\newenvironment{docspec}{\begin{quote}\noindent}{\end{quote}}% command specification environment
\newcommand{\docenv}[1]{\textsf{#1}}% environment name
\newcommand{\docpkg}[1]{\texttt{#1}}% package name
\newcommand{\doccls}[1]{\texttt{#1}}% document class name
\newcommand{\docclsopt}[1]{\texttt{#1}}% document class option name

% Package for title style
\usepackage{sectsty}
\usepackage[utf8]{inputenc}

% Sets section number style
\setcounter{secnumdepth}{3} % uncomment this, if desired
\renewcommand\thesection{\color{white}\arabic{section}}

% Sets title style
\makeatletter
  \renewcommand{\paragraph}{\@startsection{paragraph}%
    {4}{\z@}{-1ex \@plus -1ex \@minus -.3ex}%
    {0.5ex \@plus .2ex}{\normalfont\normalsize\bfseries}}
\makeatother

\makeatletter
  \renewcommand{\subsection}{\@startsection{subsection}%
    {3}{-1.8em}{-3ex \@plus -1ex \@minus -.2ex}%
    {1.5ex \@plus .2ex}
    {\hspace*{-5.5em}\fcolorbox{ltblue}{ltblue}{\parbox[c][1.0ex][b]{4em}{\phantom{space}}}
    \normalfont\large\itshape\color{ltblue}}}
\makeatother

\makeatletter
  \renewcommand{\section}{\@startsection{section}%
    {3}{-1.01em}{-3ex \@plus -1ex \@minus -.2ex}%
    {1.5ex \@plus .2ex}
    {\hspace*{-5.5em}\fcolorbox{blue}{blue}{\parbox[c][1.0ex][b]{4em}{\phantom{space}}}
    \normalfont\Large\itshape\color{blue}}}
\makeatother

% Sets theorem style
\usepackage{thmtools}
\usepackage{transparent}
\definecolor{theb}{rgb}{0.67, 0.80, 0.91}

\declaretheorem[shaded={bgcolor=Lavender,
    textwidth=30em}]{Definition} % Colorbox-styled Theorem 
\declaretheorem[shaded={bgcolor=Thistle,
    textwidth=30em}]{Theorem} % Colorbox-styled Theorem 
\declaretheorem[shaded={bgcolor=Thistle,
    textwidth=30em}]{Corollary} % Colorbox-styled Theorem
\declaretheorem[shaded={bgcolor=PeachPuff,
    textwidth=30em}]{Proposition} % Colorbox-styled Proposition
\declaretheorem[shaded={bgcolor=Thistle,
    textwidth=30em}]{Lemma} % Colorbox-styled Lemma
\declaretheorem[shaded={rulecolor=Lavender,
    rulewidth=2pt, bgcolor={rgb}{1,1,1}}]{Example-1} % Colorbounded-styled Theorem
\declaretheorem[shaded={bgcolor=PeachPuff,
    textwidth=22em}]{Remark} % Colorbox-styled Proposition
% set up pdf bookmark depth
\declaretheorem[thmbox=L]{boxtheorem L} % Theorem box L-size
\declaretheorem[thmbox=M]{Example} % Theorem box M-size
\declaretheorem[thmbox=S]{Formula} % Theorem box S-size
\declaretheorem[thmbox=S]{Statement} % Theorem box S-size

% set up pdf bookmark depth
\hypersetup{bookmarksdepth=3}

% ------------------------------------------------------------
\title{MATH255: Honours Analysis 2}
\author[Matthew He]{Matthew He}
  % if the \date{} command is left out, the current date will be used
% Beginning of the document
\begin{document}
\maketitle% this prints the handout title, author, and date

% abstract
\begin{abstract}
\noindent Abstract
\end{abstract}

\tableofcontents

% main text
\section{Metric Spaces}
\subsection{Openness and closeness}
We first gives the definition of openness and closeness of sets in metric spaces, 
and we aim to layout all their equivalent characterizations as fast as possible.

\begin{Definition}[Open and closed sets]
    Let \( (X,d) \) be a metric space and \( A \subseteq X \).
    \begin{enumerate}
        \item We say that \( A \) is open in \( (X,d) \) if for every \( x \in A\), there exists \( r > 0 \) such that \(B(x, r) \subseteq A \).
        \item We say that \( A \) is closed in \( (X, d) \) if \( X \setminus A \) is open in \( (X, d) \).
    \end{enumerate}
\end{Definition}
\marginnote{
    \begin{Remark}
        The whole space \( X \) and the empty set \( \emptyset \) are both open and closed in \( (X, d) \).
    \end{Remark}
}

\begin{Definition}[Closure, interior, and boundary]
    Let \((X, d)\) be a metric space and \( A \subseteq X \).
    \begin{enumerate}
        \item We call the closure of \( A \) in \( (X,d) \) the set \( \overline{A} \) of all 
        points in \( X \) such that \( \forall r>0 \), \( B(x, r) \cap A \neq \emptyset \).
        \item We call the interior of \( A \) in \( (X, d) \) the set \( Int(A) \) of all points in \( A \) such that
        \( \exists r > 0 \) such that \( B(x, r) \subseteq A \).
        \item We call the boundary of \( A \) in \( (X, d) \) the set \( \partial A = \overline{A} \setminus Int(A) \).
    \end{enumerate}
\end{Definition}

Some basic characterizations of the openness and closeness of sets using the definitions of closure and interior are as follows:
\begin{Theorem}[Open set and interior]
    Let \((X, d)\) be a metric space and \( A \subseteq X \). Then the following holds:
    \begin{enumerate}
        \item \( Int(A) \subseteq A \).
        \item \( Int(A) \) is open in \( (X, d) \).
        \item \(B \subset Int(A)\) for every open set \( B \subseteq A \) in \( (X, d) \).
        \item Thus, \( Int(A) \) is the largest open set in \( (X, d) \) contained in \( A \).
    \end{enumerate}
    Moreover, \( A \) is open in \( (X, d) \) iff \( Int(A) = A \).
\end{Theorem}

\begin{Theorem}[Closed set and closure]
    Let \((X, d)\) be a metric space and \( A \subseteq X \). Then the following holds:
    \begin{enumerate}
        \item \( A \subseteq \overline{A} \).
        \item \( \overline{A} \) is closed in \( (X, d) \).
        \item \( B \supseteq \overline{A} \) for every closed set \( B \supseteq A \) in \( (X, d) \).
        \item Thus, \( \overline{A} \) is the smallest closed set in \( (X, d) \) containing \( A \).
    \end{enumerate}
    Moreover, \( A \) is closed in \( (X, d) \) iff \( A = \overline{A} \).
\end{Theorem}



\paragraph{Convergence with open and close sets}
To show further characterizations of the openness and closeness,

we need to first define the limit and convergence 
of sequences in metric spaces.

\begin{Definition}[Limit and convergence]
    Let \((X, d)\) be a metric space, \( x \in X \), and \( (x_k)_{k\in \mathbb{N}} \) be
    a sequence of elements in \( X \). 
    We say that \((x_k)_k\) converges to \( x \) in \( (X, d) \) if
    and we denote \(\lim_{k\to\infty}x_k = x\) or \( x_k \to x \) as \( k \to \infty \) if
    \begin{align}
        \forall \varepsilon > 0, \; \exists k_\varepsilon \in \mathbb{N} \text{ such that } 
        \forall k \geq k_\varepsilon, d(x_k, x) < \varepsilon
    \end{align}
    Moreover, \( x \) is called the limit of the sequence \( (x_k)_k \) in \( (X, d) \).
\end{Definition}
\marginnote{
    \begin{Remark}
        Let \( (X, d) \) be a metric space and \( (x_k)_k \) be a sequence in \( X \). 
        \begin{enumerate}
            \item \(\lim_{k\to\infty} x_k = x\)in \( (X,d) \) iff \( \lim_{k\to\infty}d(x_k,x) = 0\) as a sequence in \( (\mathbb{R}, |\cdot|) \).
            \item If \(\lim_{k\to\infty} x_k = x\) in \( (X, d) \), then it's unique
            \item Every subsequence of \( (x_k)_k \) also converges to the same limit \( x \) in \( (X, d) \).
        \end{enumerate}
    \end{Remark}
}

\begin{Definition}[Cluster point (limit point)]
    \hfill \newline
    Let \( (X, d) \) be a metric space, \( A \subseteq X \), and \( x_0 \in X \).\\ 
    We say that \( x_0 \) is a cluster point of \( A \) in \( (X, d) \) if 
    \( x_0 \in \overline{A\setminus \{x_0\}} \).
\end{Definition}

\begin{Proposition}
    A point \( x_0 \) is a cluster point of a set \( A \) iff every open ball \( B(x_0, r) \) intersects \( A \setminus \{x_0\} \)
    (namely, contains a point \( x \) of \( A \) such that \( x \neq x_0 \)).
    We note that \( x_0 \) may not be in \( A \).
\end{Proposition}

\begin{Proposition}
    \hfill\newline
    \( x_0 \) is a cluster point of \( A \) iff there exists a sequence 
    \( (x_k)_{k\in \mathbb{N}} \) in \( A \) such that 
    \(\lim_{k\to \infty} x_k = x_0\) and \( x_k \neq x_0 \) for all \( k \in \mathbb{N} \).
\end{Proposition}

\begin{Definition}[Isolated point]
    Let \((X, d)\) be a metric space, \( A \subseteq X \), and \( x_0 \in A \).
    We say that \( x_0 \) is an isolated point of \( A \) in \( (X, d) \) if
    there exists \( r > 0 \) such that \( B(x_0, r) \cap A = \{x_0\} \).
\end{Definition}

\paragraph{Additional characterizations}
We can not show more characterizations of the openness and closeness of sets:
\begin{Definition}[Derived set]
    Let \((X, d)\) be a metric space, \( A \subseteq X \). The derived set 
of \( A \) in \( (X, d) \) is the set \( A' \) of all cluster points of \( A \) in \( (X, d) \).
\end{Definition}
\begin{Proposition}
    Let \((X, d)\) be a metric space and \( A \subseteq X \). Then
    the closure of \( A \) in \( (X, d) \) is the union of \( A \) and and its derived set, 
    i.e. \( \overline{A} = A \cup A' \).
\end{Proposition}

\begin{Theorem}[Close sets in metric spaces]
    For a metric space \( (X, d) \) and \( A \subseteq X \), then
    the set \( A \) is closed iff \( A \) contain all of its cluster points.
\end{Theorem}

We now characterize the closeness of sets using the convergence of sequences:
\begin{Proposition}
    Let \((X,d)\) be a metric space, \( x \in X \), and \(A \subseteq X\). Then 
    \(x\in \overline A\) iff there exists a sequence \( (x_k)_k \) in \( A \) such that \( \lim_{k\to \infty} x_k = x \).
\end{Proposition}



\begin{Theorem}[Sequential Criterion for Closedness]
    Let \((X, d)\) be a metric space and \( A \subseteq X \). Then the following statements are equivalent:
    \begin{itemize}
        \item \( A \) is closed in \( (X, d) \).
        \item For every \( x\in X \) and \((x_k)_{k\in \mathbb{N}}\) sequence of elements in \( A \), \\
        if \(\lim_{k\to \infty} x_k = x\) in \( (X, d) \), then \( x \in A \).
    \end{itemize}
\end{Theorem}


More properties of open and closed sets includes:
\begin{Proposition}[Union and intersection of open set]
    Let \((X, d)\) be a metric space,
    \begin{enumerate}
        \item The intersection of any family \( \mathcal{F} \) of open sets in \( (X, d) \) is open in \( (X, d) \).
        \item The union of any family \( \mathcal{F} \) of open sets in \( (X, d) \) is open in \( (X, d) \).
    \end{enumerate}
\end{Proposition}

\begin{Proposition}[Union and intersection of closed set]
    Let \((X, d)\) be a metric space,
    \begin{enumerate}
        \item 
    The intersection of any family \( \mathcal{F} \) of closed sets in \( (X, d) \) is closed in \( (X, d) \).
    \item The union of any finite family \( \mathcal{F} \) of closed sets in \( (X, d) \) is closed in \( (X, d) \).
    \end{enumerate}
\end{Proposition}

\begin{Proposition}[Set difference of open and closed sets]
    Let \((X, d)\) be a metric space,
    \begin{enumerate}
        \item If \( A_1 \) is an open set in \( X \) and \( A_2 \) is a closed set, then \( A_1\setminus A_2 \) is open in \( (X,d) \).
        \item If \(A_1\) is a closed set and \( A_2 \) is an open set in \( (X, d) \), then \( A_1\setminus A_2 \) is closed in \( (X, d) \).
    \end{enumerate}
\end{Proposition}

More characterizations of interior and closure:
\begin{Proposition}
    Let \( (X,d) \) be a metric space and \( A\subseteq X \)
    \begin{enumerate}
        \item \( \overline{X \setminus A} = X\setminus Int(A)\), and we have \(Int(X \setminus A) = X\setminus \overline{A}\).
        \item \( \partial A = \overline{A} \setminus Int(A) = \overline{A} \cap \overline{X\setminus A}\).
    \end{enumerate}
\end{Proposition}
Thus, boundary is the two-way limit point, it's both a limit point of \( A \) and a limit point of \( X \setminus A \).



\subsection{Completeness and compactness}

\begin{Definition}[Cauchy sequence]
    Let \((X, d)\) be a metric space and \( (x_k)_{k\in \mathbb{N}} \) be a sequence of elements in \( X \).
    We say that \((x_k)_k\) is a Cauchy sequence in \( (X, d) \) if
    \begin{align}
        \forall \varepsilon > 0, \; \exists k_\varepsilon \in \mathbb{N} \text{ such that } 
        \forall m, n \geq k_\varepsilon, d(x_m, x_n) < \varepsilon
    \end{align}
\end{Definition}

\begin{Proposition}
    Every converging sequence in \((X, d)\) is a Cauchy sequence in \( (X, d) \).
\end{Proposition}

\begin{Definition}[Completeness, Banach and Hilbert space]
    We say that a metric space \((X, d)\) is complete if every 
    Cauchy sequence in \( (X, d) \) converges to some limit in \( (X, d) \).
    \begin{itemize}
        \item A complete normed vector space is called a Banach space.
        \item A complete inner product space is called a Hilbert space.
    \end{itemize}
\end{Definition}

\begin{Proposition}
    Let \((X,d)\) be a metric space and \(Y\subseteq X\),
    \begin{enumerate}
        \item If \((Y,d)\) is complete, then \(Y\) is closed in \((X,d)\).
        \item If \((X,d)\) is complete and \( Y \) is closed in \((X,d)\), then \((Y,d)\) is complete.
    \end{enumerate}
\end{Proposition}


\begin{Definition}[Compactness]
    Let \((X, d)\) be a metric space and \( A \subseteq X \). 
    We say that \( A \) is compact (or sequentially compact) in \( (X, d) \) if for every
    sequence \( (x_k)_{k\in \mathbb{N}} \) of elements in \( A \),
    there exists \( x \in A \) and a subsequence \( (x_{k_j})_{j\in \mathbb{N}} \) such that 
    \(\lim_{j\to \infty} x_{k_j} = x\) in \( (X, d) \).

    If this is true for \(A = X\), then we say that \((X, d)\) is compact (or sequentially compact).
\end{Definition}
\marginnote{
    Compactness forces convergence and existence of limits.
}   

\begin{Proposition}
    Let \((X, d)\) be a compact metric space and \( A \subseteq X \). 
    Then \( A \) is compact in \( (X, d) \) iff \( A \) is closed in \( (X, d) \).
\end{Proposition}

\begin{Theorem}[Heine-Borel Theorem]
    Let \(X, \|\cdot\|\) be a normed vector space of finite
    dimension and \( A \subseteq X \). 
    Then, \( A \) is compact in \( (X, \|\cdot\|) \) iff \( A \) is closed and bounded in \( (X, \|\cdot\|) \).
\end{Theorem}

\begin{Proposition}[Finite subcover of compact sets]
    \hfill\newline
    Let \((X, d)\) be a metric space and \( A \subseteq X \). 
    If \( A \) is compact in \( (X, d) \), then for every collection 
    \(\{U_i\}_{i \in I}\) of open subsets of \( (X, d) \) such that \( A \subseteq \bigcup_{\alpha \in I} U_\alpha \),
    there exists a finite sub-collection \(\{U_{i_{j}}\}_{1 \leq j \leq N}\) where \( N \in \mathbb{N} \) and 
    \( i_j \in I \) for all \( 1 \leq j \leq N \) such that \( A \subseteq \bigcup_{j=1}^{N} U_{i_j} \).
\end{Proposition}


\begin{Corollary}
    Let \((X,d)\) be a metric space and \( A \subseteq X \). 
    If \( A \) is compact in \( (X, d) \), then for every \(r>0\), 
    there exists a finite number \(N_r \in \mathbb{N}\) of elements \( x_1, \cdots, x_{N_r} \in A \) such 
    that \( 
    A \subseteq \bigcup_{j=1}^{N_r} B(x_j, r)
    \).
\end{Corollary}

\begin{Corollary}
    Let \((X, \|\cdot\|)\) be a normed vector space of infinite dimension. 
    Then, the closed unit ball \( \{x \in X: \|x\| \leq 1\} \) is not compact in \( (X, \|\cdot\|) \).
    (Hence the Heine-Borel theorem does not hold in infinite dimension.)
\end{Corollary}

\begin{Proposition}
    Every normed vector space of finite dimension is Banach.
\end{Proposition}

\subsection{Mapping between metric spaces}


\begin{Definition}[Limit of function value]
    Let \( (X, d_X) \) and \( (Y, d_Y) \) be two metric spaces, \( A \subseteq X \), \( x_0\) be a 
    cluster point of \( A \) and \( y_0 \in Y \).
    \begin{enumerate}
        \item For \(f:A\to Y\), we say that \( f(x) \) converges to \( y_0 \) as \( x\to x_0 \)
         and we denote \(\lim_{x \to x_0}f(x) = y_0\) if 
         \begin{align}
            \forall \varepsilon > 0, \; \exists \delta_\varepsilon > 0:\;
            \forall x \in A, 0 < d_X(x, x_0) < \delta_\varepsilon\Rightarrow d_Y(f(x), y_0) < \varepsilon
         \end{align}
        \item When \( (Y, d_Y) = (\mathbb{R}, |\cdot|) \), for \( f:A\to \mathbb{R} \), we say that 
        \( f(x)\) tends to \( \pm\infty \) as \( x \to x_0 \) and we denote \(\lim_{x \to x_0}f(x) = \pm\infty\) if
        \begin{align}
            \forall M > 0, \; \exists \delta_M > 0:\;
            \forall x \in A, 0 < d_X(x, x_0) < \delta_M\Rightarrow f(x) > M \text{ or } f(x) < -M
        \end{align}
    \end{enumerate}
\end{Definition}







\subsection{Continuous functions on compact sets}

\section{Infinite Series}
\begin{Definition}[Partial sum]
    Let \((X, \|\cdot\|)\) be a normed vector space over 
    \(\mathbb{R}\) or \(\mathbb{C}\).
    \( k_0\in \mathbb{Z} \) and \( (a_k)_{k\geq k_0}\) by a sequence in \( X \).
    For every \( n \geq k_0 \), \(S_n = \sum_{k=k_0}^{n}a_k\) is called the partial sum of 
    order \( n \) of \( (a_k)_k \).

    If \(\lim_{n \to \infty} S_n = S\) in \( (X, \|\cdot\|) \), 
    for some \( S \in X \), then we write 
    \begin{align}
        \sum_{k=k_0}^{\infty}a_k = S
    \end{align}
    and we say that the series \( \sum_{k=k_0}^{\infty}a_k \) and that 
    \( S \) is its sum. If \((S_n)_{n\geq k_0}\) does not converge, 
    then we say that the series \( \sum_{k=k_0}^{\infty}a_k \) diverges.
\end{Definition}
\marginnote{
    \vspace{-10em}
    \begin{Remark}[Examples]
    \hfill\newline
        \begin{itemize}
            \item \textbf{Geometric series}
            \(a_k = r^k\) for some \( r \in \mathbb{R} \),
            \begin{align}
                S_n = \sum_{k=0}^{n}r^k = 
                \begin{cases}
                \dfrac{1-r^{n+1}}{1-r} & \text{if } r \neq 1\\
                n+1 & \text{if } r = 1
                \end{cases}
            \end{align}
            It converges iff \( |r| < 1 \) with \(S = \frac{1}{1-r}\) 
            since \( \lim_{n \to \infty}r^{n+1} = 0 \) when \( |r| < 1 \).
            \item \textbf{Series of matrices}
            For example, \(\sum_{k=0}^{\infty}A^k = (I-A)^{-1}\) for all 
            \( n \times n \) real value matrices such that \( \|A\|_{2,2} = \max_{\|x\|_2 = 1}\|Ax\|_2 < 1 \).
             \item \textbf{p-series}
             \(\sum_{k=1}^{\infty} \frac{1}{k^\alpha}\) converges iff \( \alpha > 1 \).
        \end{itemize}
    \end{Remark}
}

\begin{Proposition}
    \( \forall \; k_0' > k_0 \), \(\sum_{k=k_0}^{\infty}a_k\) converges iff \(\sum_{k=k_0'}^{\infty}a_k\) converges.
    Therefore, we will sometimes assume without loss of generality that \( k_0 = 0 \) or \(1 \) .
\end{Proposition}
Proof: 
\(\forall n \geq k_0'\), 
\(\sum_{k=k_0}^{n} a_k = \sum_{k=k_0'}^{n} a_k + \sum_{k=k_0}^{k_0'-1}a_k\)
where the second term is independent of \( n \).
Hence, \(\sum_{k=k_0}^{\infty}a_k\) converges iff \(\sum_{k=k_0'}^{\infty}a_k\) converges.

\hfill\qedsymbol

\begin{Proposition}
    If \(\sum_{k=k_0}^{\infty}a_k\) converges, then \( a_k \to 0 \) as \( k \to \infty \).
\end{Proposition}
Proof: If \(\sum_{k=k_0}^{\infty}a_k\) converges, then \( a_n = S_n - S_{n-1} \to S - S = 0 \) as \( n \to \infty \).
\hfill\qedsymbol

\begin{Proposition}
    If \(\sum_{k=k_0}^{\infty}a_k\) and \(\sum_{k=k_0}^{\infty}b_k\) converge, 
    then \(\forall \lambda, \mu \in \mathbb{R}\) or \( \mathbb{C} \),
    \(\sum_{k=k_0}^{\infty}(\lambda a_k + \mu b_k) \to 
    \lambda \sum_{k=k_0}^{\infty}a_k + \mu \sum_{k=k_0}^{\infty}b_k\).
\end{Proposition}
Proof follows from \(
\sum_{k=k_0}^{n}(\lambda a_k + \mu b_k) = \lambda \sum_{k=k_0}^{n}a_k + \mu \sum_{k=k_0}^{n}b_k
\).
\hfill \qedsymbol

\begin{Theorem}[Cauchy criterion for convergence]
    If \((X, \|\cdot\|)\) is Banach, then \(\sum_{k=k_0}^{\infty}a_k\) converges iff
    \(\forall \varepsilon > 0\), \(\exists n_\varepsilon \in \mathbb{N}\) such that
    \begin{align}
        \forall m > n \geq n_\varepsilon, \; \|\sum_{k=n+1}^{m}a_k\| < \varepsilon
    \end{align}
\end{Theorem}
Proof: If \((X, \|\cdot\|)\) is Banach, then \((S_n)_{n\geq k_0}\), 
\(S_n = \sum_{k=k_0}^{n}a_k\) converges iff it is Cauchy, i.e.
\(\forall \varepsilon > 0\), \(\exists n_\varepsilon \in \mathbb{N}\) such that
\begin{align}
    \forall m > n \geq n_\varepsilon, \; \|S_m - S_n\|= \|\sum_{k=n+1}^{m}a_k\|
     < \varepsilon
\end{align}
\hfill \qedsymbol

\begin{Definition}[Absolute and conditional convergence]
    \(\sum_{k=k_0}^{\infty}a_k\) is said to be absolutely convergent if 
    \(\sum_{k=k_0}^{\infty}\|a_k\|\) converges in \((\mathbb{R}, \|\cdot\|)\).
    It is said to be conditionally convergent if \(\sum_{k=k_0}^{\infty}a_k\) converges but 
    \(\sum_{k=k_0}^{\infty}\|a_k\|\) diverges.
\end{Definition}
\marginnote{
    \vspace{-8em}
    \begin{Remark}
        \hfill\newline
        \(\tilde{S}_n = \sum_{k=k_0}^{n}\|a_k\|\) is nondecreasing: \(\tilde{S}_{n+1}  \geq \tilde{S}_n\)
        for all \( n \in \mathbb{N} \). 
        
        Hence, \(\tilde{S}_n = \sum_{k=k_0}^{n}\|a_k\|\)
        converges iff \((\tilde{S}_n)_{n\geq k_0} \) is bounded. 
        Moreover, \(\sum_{k=k_0}^{\infty}\|a_k\| = \sup_{n \geq k_0} \tilde{S}_n\).
    \end{Remark}
}

\begin{Proposition}
    If \((X, \|\cdot\|)\) is Banach, then for every series 
    \(\sum_{k=k_0}^{\infty}a_k\) in \( (X, \|\cdot\|) \), 
    if \(\sum_{k=k_0}^{\infty}\|a_k\|\) converges absolutely,
    then \(\sum_{k=k_0}^{\infty}a_k\) converges and
    \begin{align}
        \|\sum_{k=n}^{\infty}a_k\| \leq \sum_{k=n}^{\infty}\|a_k\|, \quad \forall n \geq k_0
    \end{align}
\end{Proposition}

\marginnote{
\vspace{-8em}
\begin{Remark}
    \hfill\newline
    The converse is true, if \(\forall \sum a_k\) that converges absolutely,
    \(\sum a_k\) converges, then \((X, \|\cdot\|)\) is Banach.
\end{Remark}
}


\begin{Proposition}
    Let \( 
    \sum_{k=k_0}^{\infty}a_k \) be a series in a Banach space \( (X, \|\cdot\|) \).
    If \(\sum_{k=k_0}^{\infty}a_k\) converges absolutely, then for every bijective function 
    \(f: \{k \geq k_0\} \to \{k \geq k_0\}\) (permutation), 
    the series \(\sum_{k=k_0}^{\infty}a_{f(k)}\) converges and
    \begin{align}
        \sum_{k=k_0}^{\infty}a_{f(k)} = \sum_{k=k_0}^{\infty}a_k
    \end{align}
    We call the series \( \sum_{k=k_0}^{\infty}a_{f(k)} \) a rearrangement of \( \sum_{k=k_0}^{\infty}a_k \).
\end{Proposition}



\subsection{Series with nonnegative term}
\begin{Theorem}[Direct Comparison Test]
    Let \( \sum_{k=k_0}^{\infty}a_k \) and \( \sum_{k=k_0}^{\infty}b_k \) be two series in \( (\mathbb{R},\|\cdot\|) \)
    such that \( 0 \leq a_k \leq b_k \) for all \( k \geq k_0 \). Then
    \begin{enumerate}
        \item If \(\sum_{k=k_0}^{\infty}b_k\) converges, \\
        then \(\sum_{k=k_0}^{\infty}a_k\) converges and 
        \(\sum_{k=k_0}^{\infty}a_k \leq \sum_{k=k_0}^{\infty}b_k\).
        \item If \(\sum_{k=k_0}^{\infty}a_k\) diverges, then \(\sum_{k=k_0}^{\infty}b_k \) diverges.
    \end{enumerate}
\end{Theorem}

\begin{Corollary}[Limit Comparison Test]
    Let \( \sum_{k=k_0}^{\infty}a_k \) and \( \sum_{k=k_0}^{\infty}b_k \) be
    two series in \( (\mathbb{R},\|\cdot\|) \) such that 
    \( a_k > 0 \) and \( b_k > 0 \) for all \( k \geq k_0 \),
    and \( \lim_{k \to \infty}\frac{a_k}{b_k} = \ell\in (0,\infty)\).
    Then \( \sum_{k=k_0}^{\infty}a_k \) converges if and only if
    \( \sum_{k=k_0}^{\infty}b_k \) converges.
\end{Corollary}

\begin{Theorem}[Root Test]
    Let \( \sum_{k=k_0}^{\infty}a_k \) be a series in \( (\mathbb{R},\|\cdot\|) \)
    such that \( a_k \geq 0\) for all \( k \geq k_0 \).
    \begin{enumerate}
        \item If \( \limsup_{k \to \infty}\sqrt[k]{a_k} < 1 \), then \( \sum_{k=k_0}^{\infty}a_k \) converges.
        \item If \( \limsup_{k \to \infty}\sqrt[k]{a_k} > 1 \), then \( \sum_{k=k_0}^{\infty}a_k \) diverges.
    \end{enumerate}
    If \( \limsup_{k \to \infty}\sqrt[k]{a_k} = 1 \), then the test is inconclusive.
\end{Theorem}


\begin{Proposition}[Ratio Test]
    Let \( \sum_{k=k_0}^{\infty}a_k \) be a series in \( (\mathbb{R},\|\cdot\|) \)
    such that \( a_k >0 \) for all \( k \geq k_0 \).
    \begin{enumerate}
        \item If \( \lim\sup_{k \to \infty}\frac{a_{k+1}}{a_k} < 1 \), then 
        \( \sum_{k=k_0}^{\infty}a_k \) converges.
        \item If \( \lim\inf_{k \to \infty}\frac{a_{k+1}}{a_k} > 1 \), then
        \( \sum_{k=k_0}^{\infty}a_k \) diverges.
    \end{enumerate}
\end{Proposition}
Proof: (1) 
Assume that \( R=\lim\sup \frac{a_{k+1}}{a_k} <1\).
Then \(\forall \varepsilon \in (0, 1-R)\), \(\exists
n_\varepsilon \in \mathbb{N}\) such that 
\( 
\forall k \geq n_\varepsilon, \sup_{k\geq n}\frac{a_{k+1}}{a_k} < R + \varepsilon < 1
\), which gives 
for all \( k \geq n_\varepsilon \), \( a_{k+1} < (R+\varepsilon)a_k 
< (R+\varepsilon)^{k-n_\varepsilon + 1}a_{n_\varepsilon} \).
When \( \varepsilon \in (0, 1-R) \), \(0 < R+\varepsilon < 1\), so 
\( \sum_{k=n_\varepsilon}^{\infty}(R+\varepsilon)^{k-n_\varepsilon + 1}a_{n_\varepsilon} \) converges,
which by direct comparison implies that 
\( \sum_{k=k_0}^{\infty}a_k \) converges.

(2) Assume that \( R=\lim\inf \frac{a_{k+1}}{a_k} >1\).
Then \(\exists k_1 \geq k_0\) such that 
\(\forall k \geq k_1, \inf_{k\geq n}\frac{a_{k+1}}{a_k} > R - \varepsilon > 1\),
i.e. \( a_{k+1} > a_k\), which implies 
\( a_k \not\to 0\) as \( k \to \infty \), thus 
\( \sum_{k=k_0}^{\infty}a_k \) diverges.
\hfill\qedsymbol

\subsection{Power series}
Now we focus on series in \( (\mathbb{R},|\cdot|) \) with terms of different sign.
\begin{Theorem}[Alternating Series Test]
    Let \( \sum_{k=k_0}^{\infty}a_k \) be a series in \( (\mathbb{R},|\cdot|) \)
    such that \(0 \leq a_{k+1} \leq a_k\) for all \( k \geq k_0 \) and 
    \( \lim_{k \to \infty}a_k = 0 \).
    Then \( \sum_{k=k_0}^{\infty}a_k\) converges and 
    \begin{align}
        0 \leq \sum_{k=k_0}^{\infty}(-1)^{k-k_0}a_k \leq a_{k_0}
    \end{align}
\end{Theorem}

\begin{Example}
    \( \sum_{k=1}^{\infty}\frac{(-1)^{k-1}}{k}
     \) converges since \(0 \leq \frac{1}{k+1} \leq \frac{1}{k} \) for all \( k \geq 1 \) and
        \( \lim_{k \to \infty}\frac{1}{k} = 0 \).
\end{Example}
Proof: By observing that 
\begin{align}
\sum_{k=k_0}^{\infty}(-1)^{k}a_k &= \sum_{k=0}^{\infty}(-1)^{k+k_0}a_{k+k_0}\\
&= (-1)^{k_0}\sum_{k=0}^{\infty}(-1)^{k}a_{k+k_0}
\end{align}
we can assume that \( k_0 = 0 \).

For every \( n \in \mathbb{N} \),
\begin{align}
    S_{2n} &= a_0 - a_1 + a_2 - \cdots + a_{2n-2} - a_{2n-1} + a_{2n} \\
    &= (a_0 - a_1) + (a_2 - a_3) + \cdots + (a_{2n-2} - a_{2n-1}) + a_{2n} \\
    &\geq 0
\end{align}
\begin{align}
    S_{2n+1} &= a_0 - a_1 + a_2 - \cdots + a_{2n-2} - a_{2n-1} + a_{2n} - a_{2n+1} \\
    &= (a_0 - a_1) + (a_2 - a_3) + \cdots + (a_{2n-2} - a_{2n-1}) + (a_{2n} - a_{2n+1}) \\
    &\leq 0
\end{align}

\begin{Definition}[Cauchy product]
    Let \(\sum_{k=0}^{\infty}a_k\) and \(\sum_{k=0}^{\infty}b_k\) be 
    series in \((\mathbb{R}, |\cdot|)\) or \((\mathbb{C}, |\cdot|)\).

    For each \( n \in \mathbb{N} \), define 
    \begin{align}
        c_n = \sum_{k=0}^{n}a_k b_{n-k} = a_0 b_n + a_1 b_{n-1} + \cdots + a_n b_0
    \end{align}
    Then the series \( \sum_{n=0}^{\infty}c_n \) is called the Cauchy product
    of \( \sum_{k=0}^{\infty}a_k \) and \( \sum_{k=0}^{\infty}b_k \).
\end{Definition}

\begin{Proposition}
    If \(\sum_{k=0}^{\infty}a_k\) and \(\sum_{k=0}^{\infty}b_k\) are series in \((\mathbb{R}, |\cdot|)\) or \((\mathbb{C}, |\cdot|)\) such that 
    \( \sum_{k=0}^{\infty}a_k \) converges and \( \sum_{k=0}^{\infty}b_k \) converges (not necessarily absolutely).
    Then the Cauchy product \( \sum_{n=0}^{\infty}c_n \) converges
    and \( \sum_{n=0}^{\infty}c_n = (\sum_{k=0}^{\infty}a_k)(\sum_{k=0}^{\infty}b_k) \).
\end{Proposition}

\begin{Example}[Counter example of Cauchy product of two conditionally converging series]
    Consider \( a_k=b_k = \frac{(-1)^k}{\sqrt{k+1}} \).
    Then \( \sum_{k=0}^{\infty}a_k \) and \( \sum_{k=0}^{\infty}b_k \) converge by the alternating series test
    and p-series convergence.

    In this case:
    \begin{align}
        |c_n| &= |\sum_{k=0}^{n}a_k b_{n-k}| \\
    \end{align}
\end{Example}


\begin{Definition}[Power series]
    Let \( (a_k)_k \) be a sequence in \( \mathbb{R} \) or
    \( \mathbb{C} \) and \( x_0 \in \mathbb{R} \) or \( \mathbb{C} \).
    The function \( x \mapsto \sum_{k=0}^{\infty}a_k(x-x_0)^k \) is 
    defined at points where it converges, is called a power series.
\end{Definition}

\begin{Proposition}[Radius of convergence]
    Let \( \sum_{k=0}^{\infty}a_k(x-x_0)^k \) be a power series in \( \mathbb{R} \) or \( \mathbb{C} \).
    Define \(R = \frac{1}{\limsup_{k\to \infty} \sqrt[k]{|a_k|}}\) with the 
    convention that \( \frac{1}{0} = \infty \) and \( \frac{1}{\infty} = 0 \).
    Then \( R \) is called the radius of convergence of the power series, and
    \begin{align}
        &\text{if } |x-x_0| < R, \text{ then } \sum_{k=0}^{\infty}a_k(x-x_0)^k \text{ converges},\\
        &\text{if } |x-x_0| > R, \text{ then } \sum_{k=0}^{\infty}a_k(x-x_0)^k \text{ diverges}.
    \end{align}
\end{Proposition}

\begin{Proposition}
    Let \( f(x) =\sum_{k=0}^{\infty}a_k(x-x_0)^k \) be a power series in \( (\mathbb{R}, |\cdot|) \) 
    with the radius of convergence \( R > 0 \). 
    Then \( f \) is infinitely differentiable in \( (x_0 - R, x_0 + R) \) 
    
\end{Proposition}



\begin{Example}
For every \( x \in (0,2) \), we can write \( 
\ln x = \sum_{k=1}^{\infty} \frac{(-1)^{k-1}}{k}(x-1)^k \).
\end{Example}
Indeed, let \(\tilde{\ln} x = \sum_{k=1}^{\infty} \frac{(-1)^{k-1}}{k}(x-1)^k\).
We want to show that \( \tilde{\ln} x = \ln x \), i.e., 
\(\tilde{\ln} e^y = y\) where \( y = \ln x \).
Observe that 
\begin{align}
    \tilde{\ln}' (x) &= \sum_{k=1}^{\infty} (-1)^{k-1}(x-1)^{k-1} 
    = \sum_{k=1}^{\infty} (1-x)^{k-1}\\
    &= \sum_{k=0}^{\infty}(1-x)^{k} \\
    &= \frac{1}{1-(1-x)} 
    = \frac{1}{x}
\end{align}
More, this is true for all \( x \in (0,2) \) 
since the radius of convergence of \(\sum_{k=0}^{\infty}(1-x)^k\) is 1.
It follows that 
\begin{align}
    \frac{d}{dy}[\tilde{\ln} e^y - y] &= e^y\;\tilde{\ln}' e^y -1 = e^y\frac{1}{e^y} -1 = 0
\end{align}
hence \( \tilde{\ln} e^y - y = \tilde{\ln} e^0 - 0 = 0 \), 
i.e. \( \tilde{\ln} x = \ln x \).
\hfill\qedsymbol

\begin{Definition}[Taylor series]
    Let \( I \) be an interval in \( R \), \( x_0 \in I \)
    and \( f: I \to \mathbb{R} \) be infinitely differentiable
    at \( x_0 \). Then the series 
    \begin{align}
        \sum_{k=0}^{\infty}\frac{f^{(k)}(x_0)}{k!}(x-x_0)^k
    \end{align}
    is called the Taylor series of \( f \) at \( x_0 \).
\end{Definition}
\marginnote{
    We've proven that every power series 
    \( 
    f(x) = \sum_{k=0}^{\infty}a_k(x-x_0)^k
    \)
    is equal to its Taylor series at \( x_0 \in (x_0 - R, x_0 + R)\).

    However, this is not true in general where \( f \) is 
    only assumed to be infinitely differentiable at \( x_0 \). For example,
    \begin{Example}
        The function \(x \mapsto e^{-\frac{1}{x^2}}\) for \( x \neq 0 \) and \( 0 \) for \( x = 0 \)
        is infinitely differentiable in \( \mathbb{R} \)
        and such that \( f^{(k)}(0) = 0 \) for all \( k \in \mathbb{N} \).
        That means that its Taylor series at \( 0 \) is the zero series.

        Thus the function itself is not equal to its Taylor series at \( 0 \).
    \end{Example}
}

\begin{Theorem}[Taylor's theorem with mean value reminder]
    Let \( x, x_0 \in \mathbb{R} \), \( x \neq x_0 \), 
    \(I = \begin{cases}
        [x_0, x] & \text{if } x > x_0\\
        [x, x_0] & \text{if } x < x_0
    \end{cases}\), 
    \( \text{int}(I) = \begin{cases}
        (x_0, x) & \text{if } x > x_0\\
        (x, x_0) & \text{if } x < x_0
    \end{cases} \) and \( f: I \to \mathbb{R} \) be 
    n times continuously differentiable in \( I \) (i.e. 
    \( f^{(n) } \) exists and is continuous in \( I \)).
    and \( n+1 \) times differentiable in \( \text{int}(I) \).
    
    Then, \(\exists x' \in \text{int}(I) \) such that
    \begin{align}
        f(x) = \underbrace{\sum_{k=0}^{n}\frac{f^{(k)}(x_0)}{k!}(x-x_0)^k}_{\text{partial sum of Taylor series at \( x_0 \)}} 
        + \underbrace{\frac{f^{(n+1)}(x')}{(n+1)!}(x-x_0)^{n+1}}_{\text{mean value reminder}}
    \end{align}
\end{Theorem}
Define \(g(t) = f(x) = \sum_{k=0}^{n}\frac{f^{(k)}(t)}{k!}(t-t)^k \).
We want to show that 
\(g(x_0) = f(x) - \frac{1}{(n+1)!}f^{(n+1)}(x')(x-x_0)^{n+1}\) for some 
\( x' \in \text{int}(I) \).

In particular, \( g(x) = f(x) - f(x) = 0 \).

\( g \) is continuous on \( I \) and differentiable on \( \text{int}(I) \) (
    by assumption on \( f \)) and 
\begin{align}
    g'(t) &= \sum_{k=0}^{n}\frac{f^{(k+1)}(t)}{k!}(t-t)^k \\
    &= \frac{f^{(n+1)}(t)}{n!}(t-t)^n \\
    &= \frac{f^{(n+1)}(t)}{(n+1)!}(t-x_0)^{n+1} - \frac{f^{(n+1)}(t)}{(n+1)!}(t-x)^{n+1}
\end{align}
Now define 
\begin{align}
    h(t) = g(t)(x-x_0)^{n+1} - g(x_0)(x-t)^{n+1}
\end{align}
so that 
\begin{align}
    h(x) &= 0\\
    h(x_0) &= 0
\end{align}
\marginnote{
    The point here is to construct a function \( h \) such 
    that it vanishes at both \( x \) and \( x_0 \) so that we can apply Rolle's theorem to it.
}
Since \( h \) is continuous on \( I \) and differentiable on \( \text{int}(I) \) 
(by assumption on \( g \)), and 
\begin{align}
 h'(t) &= g'(t)(x-x_0)^{n+1}  + (n+1)g(x_0)(x-t)^n\\
\end{align}
By Rolle's Theorem, it follows that 
\begin{align}
    \exists x' \in \text{int}(I) \text{ such that } h'(x') = 0
\end{align}

\clearpage
\section{Riemann-Stieltjes integral}
Riemann-Stieltjes integral generalizes the Riemann integral by introducing a
weighting function \( \alpha \). The standard width in the Riemann integral 
\(\Delta x_i = x_i - x_{i-1}\) is replaced by \( \Delta \alpha_i = \alpha(x_i) - \alpha(x_{i-1}) \).

This has lots of (physical) implications. For example, \(\alpha\) can be a 
cumulative distribution function of a random variable, and then the Riemann-Stieltjes integral of a function \( f \) with respect to \( \alpha \)
becomes the expectation.

In fact, it is a reformulation of what we usually write as
\( \int f(x) \rho(x) dx \) in physics or probability.

We first gives definition of the Riemann-Stieltjes integral,
then we clarify the keep intuition behind it.
\begin{Definition}[Partition]
    Let \( a, b \in \mathbb{R} \) such that \( a < b \).
    We call partition of \([a,b]\) a finite set 
    \( \{x_0, x_1, \ldots, x_n\} \), where \( n \in \mathbb{N} \), 
    of elements in \([a,b]\) such that \( a = x_0 \leq x_1 \leq \cdots \leq x_n = b \).
\end{Definition}

\begin{Definition}[Riemann-Stieltjes integrals]
    Let \( \alpha, f: [a,b] \to \mathbb{R} \) be such that \( \alpha \) 
    is increasing on \([a,b]\) and \( f \) is bounded on \([a,b]\).
    For every partition \( P = (x_0, x_1, \ldots, x_n) \) of \([a,b]\), we define 
    \begin{align}
        L(P, f, \alpha) = \sum_{i=1}^{n}m_i\Delta \alpha_i,\quad 
        U(P, f, \alpha) = \sum_{i=1}^{n}M_i\Delta \alpha_i
    \end{align}
    where \( \Delta \alpha_i  = \alpha(x_i) - \alpha(x_{i-1}) \),
    \(m_i = \inf_{x \in [x_{i-1}, x_i]}f(x)\) and \( M_i = \sup_{x \in [x_{i-1}, x_i]}f(x) \).
    We then define 
    \begin{align}
        \underline{\int_{a}^{b}}fd\alpha = \sup_{P}L(P, f, \alpha),\quad 
        \overline{\int_{a}^{b}}fd\alpha = \inf_{P}U(P, f, \alpha)
    \end{align}
    they are called the \textbf{lower and upper Riemann-Stieltjes integrals} of \( f \) with respect to \( \alpha \) on \([a,b]\) respectively.

    If \( \underline{\int_{a}^{b}}fd\alpha = \overline{\int_{a}^{b}}fd\alpha \), then 
    we denote \(\int_{a}^{b}fd\alpha = \underline{\int_{a}^{b}}fd\alpha = \overline{\int_{a}^{b}}fd\alpha \)
    and we say that \( f \) is \textbf{Riemann-Stieltjes integrable} with respect to \( \alpha \) on \([a,b]\),
    and \( \int_{a}^{b}fd\alpha \) is called the Riemann-Stieltjes integral of \( f \) with respect to \( \alpha \) on \([a,b]\).

    We denote by \( \mathcal{R}([a,b], \alpha) \) the set of all bounded functions
    \( f: [a,b] \to \mathbb{R} \) which are Riemann-Stieltjes integrable with respect to \( \alpha \) on \([a,b]\).

    In the case when \(\alpha(x) = x\), we say Riemann integral instead of Riemann-Stieltjes integral 
    and we write \( \int_{a}^{b}f = \int_{a}^{b}fd\alpha \),
    \( L(P,f) = L(P,f,\alpha) \), \(\mathcal{R}([a,b]) = \mathcal{R}([a,b], \alpha) \) and so on.
\end{Definition}
\marginnote{
\begin{Remark}
    \( \alpha \) is bounded on \([a,b] \) since \( \alpha(a) \leq \alpha(x) \leq \alpha(b)\; \forall x \in [a,b] \) and \( \alpha(a), \alpha(b) \in \mathbb{R} \).
\end{Remark}
}
We note that we have:
\begin{align}
    \left(
        \inf_{a \leq x \leq b} f(x)
    \right)
    \left(
        \alpha(b) - \alpha(a)
    \right)
    &=
    \left(
        \inf_{a \leq x \leq b} f(x)
    \right)\sum_{i=1}^{m}\Delta \alpha_i\\
    &\leq L(P, f, \alpha) \leq U(P, f, \alpha)\\
    &\leq 
    \left(\sup_{a \leq x \leq b} f(x)
    \right)\left(
        \alpha(b) - \alpha(a)
    \right)
\end{align}
since \(\inf_{a \leq x \leq b} f(x) \leq m_i \leq M_i \leq \sup_{a \leq x \leq b} f(x) \) and 
\( \Delta \alpha_i \geq 0 \), hence, 
\begin{align}
    \left(
        \inf_{a \leq x \leq b} f(x)
    \right)
    \left(
        \alpha(b) - \alpha(a)
    \right)
    &\leq \underline{\int_{a}^{b}}fd\alpha 
    \leq \overline{\int_{a}^{b}}fd\alpha 
    \leq 
    \left(\sup_{a \leq x \leq b} f(x)
    \right)\left(
        \alpha(b) - \alpha(a)
    \right)
\end{align}

Unlike the Riemann integral, the Riemann-Stieltjes integral 
is constructed through sets argument with inf/sup instead of sequences argument with limits.
The core of the construction is built on the sets of partitions, which then gives
sets of lower and upper sums. 
Instead of taking limits of \(n \to \infty\) in the Riemann integral, we 
take the supremum and infimum over the sets of partitions in the Riemann-Stieltjes integral.

Here, we highlight that taking inf/sup works like taking a generalized limit (but on a set).
(Definition of Inf/sup and the process of taking them naturally inherits (wraps/encapsulates) 
the epsilon-delta style of arguments in the sequential limit).
Moreover, these set arguments are built on the technical tools of refinement.
Refinement is like moving forward toward the direction of limit, 
(increasing \( n \) when \(n\to \infty\)). 
We will see how it works in the proof of the following propositions and the 
integrability criterion.
\begin{Definition}[Refinement]
    We say that a partition \(P^* \) of \( [a,b] \) is a refinement 
    of a partition \( P \) of \( [a,b] \) if \( P^* \supseteq P \).
    Given two partitions \( P_1 \) and \( P_2 \) of \( [a,b] \), 
    we say that  \(P_1 \cup P_2 \) is a common refinement of \( P_1 \) and \( P_2 \).
\end{Definition}
The inclusion \( P^* \supseteq P \) works like the inequality \( N \geq N_0 \) in the sequential limit.
Taking the common refinement of two partitions is like taking the \( N = \max(N_1,N_2) \)
in the epsilon-delta style of arguments in the sequential limit.

\begin{Proposition}
    If \( P^* \) is a refinement of \( P \), then 
    \begin{align}
        L(P,f,\alpha) \leq L(P^*, f, \alpha) \text{ and } U(P^*, f, \alpha) \leq U(P, f, \alpha)
    \end{align}
\end{Proposition}
Proof: We prove the results in the case where \( P^* \) contain one more point 
\( x^* \) than \( P \) (the general case then follows by induction).

Assume that \( x^* \in (x_{i-1}, x_i) \) where \( (x_{i-1}, x_i) \) are
consecutive points of \( P \). Then
\begin{align}
    L(P^*, f, \alpha) - L(P, f, \alpha) 
    &= 
    m_i^* \Delta \alpha_i^* + m_{i+1}^* \Delta \alpha_{i+1}^* - m_i \Delta \alpha_i
\end{align}
where \( m_i^* = \inf_{x \in [x_{i-1}, x^*]}f(x) \), \( m_{i+1}^* = \inf_{x \in [x^*, x_i]}f(x) \),
\( \Delta \alpha_i^* = \alpha(x^*) - \alpha(x_{i-1}) \) and \( \Delta \alpha_{i+1}^* = \alpha(x_i) - \alpha(x^*) \).

Since \([x_{i-1}, x^*] \subseteq [x_{i-1}, x_i] \) and \( [x^*, x_i] \subseteq [x_{i-1}, x_i] \), we have
\begin{align}
    m_i \leq m_i^*, \; m_i \leq m_{i+1}^*
\end{align}
Moreover, since \( \alpha \) is increasing, we have \( 
\Delta \alpha_i^*, \Delta \alpha_{i+1}^* \geq 0 \) (so \( m_i \Delta \alpha_i^* \leq m_i^* \Delta \alpha_i^* \)).

Therefore, we obtain 
\begin{align}
    L(P^*, f, \alpha) - L(P, f, \alpha) 
    &\geq m_i(\Delta \alpha_i^* + \Delta \alpha_{i+1}^* - \Delta \alpha_i)=0
\end{align}
Thus, \( L(P^*, f, \alpha) \geq L(P, f, \alpha) \), and similarly, we can obtain 
\( U(P^*, f, \alpha) \leq U(P, f, \alpha) \).
\hfill\qedsymbol


\begin{Proposition}
    \( 
        \underline{\int_{a}^{b}}fd\alpha \leq \overline{
            \int_a^b
        }fd\alpha 
    \)
\end{Proposition}
Proof: \( \forall P_1, P_2 \) partitions of \([a,b]\), let \(P^* =  P_1 \cup P_2 \) be a common
refinement of \( P_1 \) and \( P_2 \). Then by the previous proposition,
\begin{align}
    L(P_1, f, \alpha) \leq L(P^*, f, \alpha) \leq U(P^*, f, \alpha) \leq U(P_2, f, \alpha)
\end{align}
where the middle inequality follows from infimum is less than or equal to supremum.

Hence, \(
        \underline{\int_{a}^{b}}fd\alpha \leq \overline{
            \int_a^b
        }fd\alpha \) 
        follows by taking 
the supremum over \( P_1 \) and the infimum over \( P_2 \).

\marginnote{
    The idea here is like taking \(N = \max(N_1, N_2)\) in the sequential limit,
    where we can find a common ansestor \( N \) at which an inequality works for both \( N_1 \) and \( N_2 \).
}
(Here, taking the inf/sup is like taking the limit.)

\hfill\qedsymbol

Indeed, the refinement gives control on the approximation of the lower and upper sums to the integral.
In the above proposition, we show the gap between the lower and upper integrals is always exists by inserting common refinement
in the middle.
In turn, 
we can also bound the error in the approximation (gap in refinement) 
by controlling the outside gap between the lower and upper sums on previous partitions.

\begin{Proposition}
    \( f \in \mathcal{R}([a,b], \alpha) \) iff \( 
    \forall \varepsilon>0 \), \(\exists P_\varepsilon\) (partition of 
    \([a,b]\)) such that \( U(P_\varepsilon, f, \alpha) - L(P_\varepsilon, f, \alpha) < \varepsilon \).
\end{Proposition}
Proof: (\(\Rightarrow\)) Assume that \( f \in \mathcal{R}([a,b], \alpha) \).
Then by properties of \( inf \) and \( sup \), \(\forall \varepsilon>0\), 
\(\exists P_\varepsilon, P'_\varepsilon\) partition 
of \([a,b]\) such that
\begin{align}
    \underline{\int_{a}^b} f(x)d\alpha  - \frac12\varepsilon < L(P_\varepsilon, f, \alpha) 
    = \underline{\int_{a}^b} f(x)d\alpha \leq \overline{\int_{a}^b} f(x)d\alpha
    \leq U(P_\varepsilon, f, \alpha) < \overline{\int_{a}^b} f(x)d\alpha + \frac12\varepsilon
\end{align}
Let \( P_\varepsilon^* = P_\varepsilon \cup P'_\varepsilon \).
Then since \( P_\varepsilon^* \supset P_\varepsilon \) 
and \( P_\varepsilon^* \supset P'_\varepsilon \),
\begin{align}
    U(P_\varepsilon^*, f, \alpha) - L(P_\varepsilon^*, f, \alpha)
    &\leq U(P_\varepsilon', f, \alpha) - L(P_\varepsilon, f, \alpha)\\
    &< (\overline{\int_{a}^b} f(x) d\alpha + \frac{1}{2}\varepsilon) - (\underline{\int_{a}^b} f(x)d\alpha - \frac{1}{2}\varepsilon)\\
    &= \overline{\int_{a}^b} f(x) d\alpha - \underline{\int_{a}^b} f(x)d\alpha + \varepsilon\\
    &= \varepsilon
\end{align}
where \( 
\overline{\int_{a}^b} f(x) d\alpha - \underline{\int_{a}^b} f(x)d\alpha =0
 \)
since \( f \in \mathcal{R}([a,b], \alpha) \).
\marginnote{
    Thus, the gap on a further refinement (like larger \( n \) in sequential proof) is less than the 
    gap on the previous partitions.
}

\noindent (\(\Leftarrow\)) 
Assume that \(\forall \varepsilon>0\), 
\(\exists P_\varepsilon\) partition of \([a,b]\) 
such that \( U(P_\varepsilon, f, \alpha) - L(P_\varepsilon, f, \alpha) < \varepsilon \).
We have 
\begin{align}
    L(P_\varepsilon, f, \alpha) \leq \underline{\int_{a}^b} f(x)d\alpha
    \leq \overline{\int_{a}^b} f(x)d\alpha \leq U(P_\varepsilon, f, \alpha)
\end{align}
It follows that 
\begin{align}
    0 \leq \overline{\int_{a}^b} f(x)d\alpha - \underline{\int_{a}^b} f(x)d\alpha 
    \leq U(P_\varepsilon, f, \alpha) - L(P_\varepsilon, f, \alpha) \leq \varepsilon
\end{align}
Letting \( \varepsilon \to 0 \), we have
\begin{align}
    \overline{\int_{a}^b} f(x)d\alpha - \underline{\int_{a}^b} f(x)d\alpha = 0
\end{align}
therefore \( f \in \mathcal{R}([a,b], \alpha) \).
\hfill\qedsymbol
\marginnote{
    Look at Question 2 in Assignment 6.
}

\begin{Theorem}[Riemann-Stieltjes integrability]
    Let \( f, \alpha [a,b] \to \mathbb{R} \) 
    be such that \( f \) is bounded in \([a,b]\)
    and \( \alpha \) increasing on \([a,b]\).

    Let \( D \) be the set of points in \([a,b]\) where \( f \) is
    not continuous. Assume that \( D \) is countable
    and that \( \alpha \) is continuous at every point in \( D \).
    Then \( f \in \mathcal{R}([a,b], \alpha) \).
\end{Theorem}
\marginnote{
    Continuous functions are always Riemann-Stieltjes integrable.
    If the function is not continuous, if 
    the set of discontinuity points is countable, and 
    if \( \alpha \) is continuous at every discontinuity point, then the function is also Riemann-Stieltjes integrable.
}
Let \(\varepsilon>0\). For every \( x \in [a,b]\setminus D\), 
since \( f \) is continuous at \( x \), there exists 
\(\delta_x > 0 \) such that \( |f(y) - f(x)| < c\varepsilon \) for all \( y \in (x-\delta_x, x+\delta_x) \cap [a,b] \),
where the constant \(c>0\) will be determined later.

Since \( D \) is countable, \(\exists (x_k)_k \) in \([a,b]\) 
such that \(D = \{x_k:k\in \mathbb{N}\}\).
Since \( \alpha \) is continuous at \(x_k\), 
there exists \( \delta_k > 0 \) such that 
\( |\alpha(y) - \alpha(x_k)| < c_k\varepsilon \) 
for all \( y \in (x_k - \delta_k, x_k + \delta_k) 
\cap [a,b] \),
where the constant \( c_k > 0 \) will be determined later.

Since \( [a,b] \) is compact, 
\( 
    [a,b] \subset \bigcup_{x \in [a,b]} (x-\delta_x, x+\delta_x)
\),
\marginnote{
    \(\delta_x\) is defined for every \( x \in [a,b] \), 
    since either it's defined by the continuity of \( f \) at \( x \) (if \( x \in [a,b]\setminus D \)) or by the continuity of \( \alpha \) at \( x \) (if \( x \in D \)).
}
there exists a finite subcover, i.e. \(\exists y_1 < \cdots < y_{N_\varepsilon} \in [a,b] \)
such that 
\begin{align}
    [a,b] \subset \bigcup_{i=1}^{N_\varepsilon}(y_i - \delta_{y_i}, y_i + \delta_{y_i})
\end{align}
By ordering the points 
\(a,b,y_i-\delta_{\varepsilon, y} (\text{ if } \geq a)\), 
and \(y_i+\delta_{\varepsilon, y} (\text{ if } \leq b)\),
we then obtain a partition \( P_\varepsilon = \{x_0, \cdots, x_{m_\varepsilon}\} \)
such that  \(\forall j \in \{0, \cdots, m_\varepsilon\}\),
\(\exists i_j \in \{1, \cdots, N_\varepsilon\}\) such 
that \( (x_{j-1}, x_j) \subset (y_{i_j} - \delta_{y_{i_j}}, y_{i_j} + \delta_{y_{i_j}})\cap [a,b] \).
\marginnote{
    Important trick here. And \(i_j\) is not a subsequence here, but a function that maps 
    the paritition index \( j \) to the index of the subcover \( i_j \).
}

We again separate this into two cases (depending on whether \( y_{i_j} \in D \) or not):
\begin{itemize}
    \item If \( y_{i_j} \not\in D\), then 
    \begin{align}
    M_j - m_j &= \sup_{x \in [x_{j-1}, x_j]}f(x) - \inf_{x \in [x_{j-1}, x_j]}f(x)\\
    &\leq \left(
    \sup_{x \in [x_{j-1}, x_j]}f(x) - f(y_{i_j})\right) + \left(f(y_{i_j}) - \inf_{x \in [x_{j-1}, x_j]}f(x)\right)\\
    &< 2c\varepsilon
    \end{align}
    \item If \( y_{i_j} \in D \), then we have
    \begin{align}
    M_j - m_j \leq M-m &= \sup_{x \in [a,b]}f(x) - \inf_{x \in [a,b]}f(x) 
    \\
    \Delta \alpha_j = \alpha(x_j) - \alpha(x_{j-1}) &\leq |\alpha(x_j) - \alpha(y_{i_j})| + |\alpha(y_{i_j}) - \alpha(x_{j-1})|\\
     &< 2c_k\varepsilon
    \end{align}
\end{itemize}
\marginnote{
    In the discontinuous regime, we have no control on the gap \( M_j - m_j \),
    however, we can use the continuity of \( \alpha \) to control the gap \( \Delta \alpha_j \).
}
Therefore, \begin{align}
    U(P_\varepsilon, f, \alpha) - L(P_\varepsilon, f, \alpha)
    &= \sum_{j=1}^{m_\varepsilon}(M_j - m_j)\Delta \alpha_j\\
    &= \sum_{j=1, y_{i_j} \in D}^{m_\varepsilon}(M_j - m_j)\Delta \alpha_j + \sum_{j=1, y_{i_j} \not\in D}^{m_\varepsilon}(M_j - m_j)\Delta \alpha_j\\
    &\leq 2c\varepsilon \underbrace{\sum_{j=1, y_{i_j} \not\in D}^{m_\varepsilon}\Delta \alpha_j}_{\leq \sum_{j=1}^{m_\varepsilon}\Delta \alpha_j = \alpha(b) - \alpha(a)} + 
    (M-m) \underbrace{ \sum_{j=1, y_{i_j} \in D}^{m_\varepsilon}\Delta \alpha_j }_{\leq 2\varepsilon\sum_{k=1}^{\infty}c_k}
\end{align}
the second inequality comes from the fact that \( D = \{d_k : k\in \mathbb{N}\} \) and \(\Delta \alpha_j < 2c_k\varepsilon \) when \( y_{i_j} = d_k \).
\marginnote{
    (Ask in Office Hour)
}

Therefore, if we choose \( c = \frac{1}{4(\alpha(b) - \alpha(a))} \) and \( c_k = \frac{1}{4(M-m)2^k} \) (if \( m < M \)),
so that \( 2(M-m)\sum_{k=1}^{\infty}c_k = \frac{1}{2} \),
then we have
\begin{align}
    U(P_\varepsilon, f, \alpha) - L(P_\varepsilon, f, \alpha) < \frac{1}{2}\varepsilon + \frac{1}{2}\varepsilon = \varepsilon
\end{align}
Since this is true for all \( \varepsilon > 0 \), it follows that \( f \in \mathcal{R}([a,b], \alpha) \).
\hfill\qedsymbol
One of the critical part of the proof is to create a shrinking \(c_k\) to control error on infinitely many 
discontinuity points.

\begin{Proposition}
    Let \( I \) be an interval in \( \mathbb{R} \) and \( f: I \to \mathbb{R} \) be 
    monotone on \( I \) and \( D \) be the set of points at which 
    \( f \) is not continuous, then \( D \) is countable.
\end{Proposition}
\marginnote{
    This shows that the countable discontinuity condition 
    in the previous theorem is automatically satisfied for monotone functions.
    \begin{Remark}
        If \( \alpha \) is continuous, then every monotone functino is Riemann-Stieltjes integrable with respect to \( \alpha \).
    \end{Remark}
}
Proof: Assuming that \( f \) is increasing (otherwise
we can consider \( -f \)). For each \( x \in D \),
we have that \( \lim_{t\to x^-}f(t) \) and \( \lim_{t\to x^+}f(t) \) exist 
since \( f \) is increasing. Moreover, 
\begin{align}
    \lim_{t\to x^-}f(t)  < \lim_{t\to x^+}f(t)
\end{align}
since \( x\in D \). By density of \( \mathbb{Q} \),
it follows that \( \exists q_x \in \mathbb{Q} \)
such that 
\begin{align}
    \lim_{t\to x^-}f(t) < q_x < \lim_{t\to x^+}f(t)
\end{align}
Moreover, since \( f \) is increasing, if \( x < y \),
then 
\begin{align}
    q_x \leq \lim_{t\to x^+}f(t) \leq \lim_{t\to y^-}f(t) < q_y
\end{align}
Therefore, the map (\( D \to \mathbb{Q} \)) \( x\mapsto q_x \) is 
injective. Since \( \mathbb{Q} \) is countable, \( D \) is countable.
\hfill\qedsymbol

\begin{Proposition}[Integrability of composition]
    Let \( a,b,c,d \in \mathbb{R} \) be such that \( a<b, c<d \), 
    \( \alpha: [a,b] \to \mathbb{R} \) be increasing, and \( f: [a,b] \to [c,d] \)
    and \( \Phi: [c,d] \to \mathbb{R} \). If 
    \( f \in \mathcal{R}([a,b], \alpha) \) and \( \Phi \) is continuous, then
    \( \Phi \circ f \in \mathcal{R}([a,b], \alpha) \).
\end{Proposition}
\marginnote{
    Implication: if we have a function \( f \) that is Riemann-Stieltjes integrable, 
    is the absolute value, or power of \( f \) also Riemann-Stieltjes integrable?
}
Proof: Since \( [c,d] \) is compact and \( \Phi \)
continuous on \( [c,d] \), we obtain that 
\( \Phi \) is uniformly continuous on \( [c,d] \). 
Thus, \(\forall \varepsilon > 0\), \(\exists \delta_\varepsilon > 0\) 
such that \(\forall x,y\in [c,d]\), 
\begin{align}
    |x-y| < \delta_\varepsilon \implies |\Phi(x) - \Phi(y)| < \lambda\varepsilon
\end{align}
where the constant \( \lambda > 0 \) will be determined later.
\marginnote{
    The trick is to keep a constant \( \lambda \) like the proof before.
}

On the other hand,  since \( f \in \mathcal{R}([a,b], \alpha) \),
\(\exists P_\varepsilon \) partition of \([a,b]\) such that
\begin{align}
    U(P_\varepsilon, f, \alpha) - L(P_\varepsilon, f, \alpha) 
    = \sum_{i=1}^{n} (M_i - m_i) \Delta \alpha_i < \lambda_\varepsilon
\end{align}
where the constant \( \lambda_\varepsilon > 0 \) will be determined later.
\marginnote{
    Recall that \( M_i = \sup_{x \in [x_{i-1}, x_i]}f(x) \) and \( m_i = \inf_{x \in [x_{i-1}, x_i]}f(x) \).
}
We want to show that 
\begin{align}
 U(P_\varepsilon, \Phi \circ f, \alpha) - L(P_\varepsilon, \Phi \circ f, \alpha) 
 = \sum_{i=1}^{n} (M'_i - m'_i) \Delta \alpha_i < \varepsilon
\end{align}
and we need to estimate the term product \( (M'_i - m'_i) \Delta \alpha_i \) for each \( i \in \{1, \cdots, n\} \).

For all \( i \in \{1, \cdots, n\} \), we have
\begin{align}
    M_i' - m_i' \leq \begin{cases}
    \lambda\varepsilon & \text{if } M_i - m_i < \delta_\varepsilon\\
    M-m & \text{if } M_i - m_i \geq \delta_\varepsilon
    \end{cases}
\end{align}
\marginnote{
    In the first case, as soon as we have one \( \varepsilon \) in the product term, it 
    will be small.
    In the second case, we can only say that \( M_i' - m_i' \) 
    is bounded, but fortunately, we have control on the other term \( \Delta \alpha_i \).
}
where \( M = \sup_{x \in [c,d]}\Phi(x) \) and \( m = \inf_{x \in [c,d]}\Phi(x) \).
This gives:
\begin{align}
    \sum_{i=1}^{n}(M_i' - m_i')\Delta \alpha_i 
    \leq \underbrace{\lambda\varepsilon \sum_{i=1,M_i - m_i < \delta_\varepsilon}^{n}\Delta \alpha_i}_{
        \leq \lambda\varepsilon (\alpha(b) - \alpha(a))
    }
    + (M-m)\sum_{i=1, M_i - m_i \geq \delta_\varepsilon}^{n}\Delta \alpha_i
\end{align}
The first term is easy to control, the second term is tricker.
But we haven't use the Riemann-Stieltjes integrability of \( f \) yet.
We have:
\begin{align}
    (M-m)\sum_{i=1, M_i - m_i \geq \delta_\varepsilon}^{n}\Delta \alpha_i
    &\leq (M-m)\sum_{i=1, M_i - m_i \geq \delta_\varepsilon}^{n}\frac{M_i - m_i}{\delta_\varepsilon} \Delta \alpha_i
\end{align}
Since \(  M_i - m_i \geq \delta_\varepsilon\).
\marginnote{
    We introduce the term \( \frac{M_i - m_i}{\delta_\varepsilon} \) so that we can use the 
    fact that \( \sum_{i=1}^{n}(M_i - m_i)\Delta \alpha_i < \lambda_\varepsilon \)
    from the integrability.
}
Since \( \sum_{i=1}^{n}(M_i - m_i)\Delta \alpha_i < \lambda_\varepsilon \), we have
\begin{align}
    (M-m)\sum_{i=1, M_i - m_i \geq \delta_\varepsilon}^{n}\Delta \alpha_i
    &\leq \frac{M-m}{\delta_\varepsilon} \sum_{i=1}^{n}(M_i - m_i)\Delta \alpha_i < \frac{M-m}{\delta_\varepsilon} \lambda_\varepsilon
\end{align}

Now, choosing \( \lambda = \frac{1}{2(\alpha(b) - \alpha(a))} \) and \( \lambda_\varepsilon = \frac{\delta_\varepsilon\varepsilon}{2(M-m)} \), we have
\begin{align}
    \sum_{i=1}^{n}(M_i' - m_i')\Delta \alpha_i \leq \frac{1}{2}\varepsilon + \frac{1}{2}\varepsilon = \varepsilon
\end{align}
\hfill \qedsymbol

\begin{Proposition}
    Let \( a, b \in \mathbb{R} \) be such that \( a<b \)
    and \( \alpha: [a,b]\to \mathbb{R} \).
    \begin{enumerate}
        \item \(\forall f,g \in \mathcal{R}([a,b], \alpha)\), \( f+g \in \mathcal{R}([a,b], \alpha) \)\\
        and 
        \begin{align}
            \int_{a}^b (f+g)d\alpha = \int_{a}^b fd\alpha + \int_{a}^b gd\alpha
        \end{align}
        \item \(\forall \lambda \in \mathbb{R}, f \in \mathcal{R}([a,b], \alpha)\),
        we have \( \lambda f \in \mathcal{R}([a,b], \alpha) \) and\\
        \begin{align}
            \int_{a}^b \lambda f d\alpha = \lambda \int_{a}^b fd\alpha
        \end{align}
        \item \(\forall f,g \in \mathcal{R}([a,b], \alpha)\) such that \( f(x) \leq g(x) \) on \( [a,b] \)
        implies that 
        \begin{align}
            \int_{a}^b fd\alpha \leq \int_{a}^b gd\alpha
        \end{align}
        \item \(\forall f \in \mathcal{R}([a,b], \alpha)\), \( |f| \in \mathcal{R}([a,b], \alpha) \) and
        \begin{align}
            \left|\int_{a}^b fd\alpha\right| \leq \int_{a}^b |f|d\alpha
        \end{align}
        \item \(\forall f,g \in \mathcal{R}([a,b], \alpha)\),
        \( fg \in \mathcal{R}([a,b], \alpha) \) (but we don't have a simple formula for \( \int_{a}^b fg d\alpha \)).
        \item \(\forall f \in \mathcal{R}([a,b], \alpha)\), \( c \in (a,b) \) 
        we have 
        \begin{align}
            \int_{a}^b fd\alpha = \int_{a}^c fd\alpha + \int_{c}^b fd\alpha
        \end{align}
    \end{enumerate}
\end{Proposition}

Proof: 
(1.)
\( \forall P, P' \) partition of \([a,b]\), letting 
\( P^* = P \cup P' \), we have
\begin{align}
    \overline{\int_{a}^{b}}(f+g)d\alpha &\leq U(P^*, f+g, \alpha)\\
    \intertext{since \( \sup(f+g)\leq \sup f + \sup g \)}
    &\leq U(P^*, f, \alpha) + U(P^*, g, \alpha)
\end{align}
By taking the infimum over \( P \) and \( P' \), we have
\begin{align}
    \overline{\int_{a}^{b}}(f+g)d\alpha &\leq \overline{\int_{a}^{b}}fd\alpha + \overline{\int_{a}^{b}}gd\alpha\\
    &= \int_{a}^{b} fd\alpha + \int_{a}^{b} gd\alpha \qquad (f,g \in \mathcal{R}([a,b], \alpha))
\end{align}
Similarly, we can show that
\begin{align}
    \underline{\int_{a}^{b}}(f+g)d\alpha &\geq \int_{a}^{b} fd\alpha + \int_{a}^{b} gd\alpha
\end{align}
Therefore, 
\begin{align}
    \int_{a}^{b}f d\alpha + \int_{a}^{b} gd\alpha \leq \underline{\int_{a}^{b}}(f+g)d\alpha \leq \overline{\int_{a}^{b}}(f+g)d\alpha \leq \int_{a}^{b} fd\alpha + \int_{a}^{b} gd\alpha
\end{align}
hence \( f+g \in \mathcal{R}([a,b], \alpha) \) and
\begin{align}
    \int_{a}^{b}(f+g)d\alpha = \int_{a}^{b} fd\alpha + \int_{a}^{b} gd\alpha
\end{align}

(2.) follows from 
\begin{align}
        \sup(\lambda f) = \lambda \sup f \text{ and } \inf(\lambda f) 
        = \lambda \inf f & \text{if } \lambda \geq 0
\end{align}
and \( \sup(-f) = -\inf(f) \), \( \inf(-f) = -\sup(f) \).
This gives 
\begin{align}
    U(P, -f, \alpha) = -L(P, f, \alpha) \text{ and } L(P, -f, \alpha) = -U(P, f, \alpha)\\
    \overline{\int_{a}^{b}}(-f)d\alpha = -\underline{\int_{a}^{b}}fd\alpha 
    \text{ and } \underline{\int_{a}^{b}}(-f)d\alpha = -\overline{\int_{a}^{b}}fd\alpha
\end{align}

(3.) follows from \( f\leq g \implies \begin{cases}
    \sup f \leq \sup g\\
    \inf f \leq \inf g
\end{cases} \).

(4.) follows from absolute function is continuous, so
\( |f| \) is Riemann-Stieltjes integrable. 
We apply (2. \& 3.) with \(-|f| \leq f \leq |f| \) to obtain the desired inequality.
\marginnote{We have to use two side inequality to bound absolute value.}

(5.) Obverse that \( fg = \frac{1}{2}((f+g)^2 - f^2 - g^2) \). 
Using (1, 2) and the fact that \(x\mapsto x^2\) is continuous.

(6.) Assignment 6.
\hfill \qedsymbol



\begin{Proposition}
    Let \( a,b \in \mathbb{R} \) such that \( a<b \)
    and \(\alpha, \alpha_1, \alpha_2: [a,b] \to \mathbb{R} \)
    be increasing on \([a,b]\). Then:
    \begin{enumerate}
        \item \( \forall f \in \mathcal{R}([a,b], \alpha_1) \cap \mathcal{R}([a,b], \alpha_2) \),
        \( f \in \mathcal{R}([a,b], \alpha_1 + \alpha_2) \) and
        \begin{align}
            \int_{a}^b fd(\alpha_1 + \alpha_2) = \int_{a}^b fd\alpha_1 + \int_{a}^b fd\alpha_2
        \end{align}
        \item \( \forall f \in \mathcal{R}([a,b], \alpha) \), \(\lambda \in (0,\infty)\), 
        \( f \in \mathcal{R}([a,b], \lambda \alpha) \) and
        \begin{align}
            \int_{a}^b fd(\lambda \alpha) = \lambda \int_{a}^b fd\alpha
        \end{align}
    \end{enumerate}
\end{Proposition}


\begin{Theorem}
    Let \( a,b \in \mathbb{R} \) be such that \( a<b \) and 
    \(f\in \mathcal{R}([a,b])\). Then the function \( F: [a,b] \to \mathbb{R} \) defined 
    by \( F(x) = \int_{a}^x f , \; \forall x \in [a,b] \) is 
    lipchitz continuous on \([a,b]\). Moreover, 
    if \( f \) is continuous at \( x_0 \in [a,b] \), then \( F \) is differentiable at \( x_0 \) and
    \begin{align}
        F'(x_0) = f(x_0)
    \end{align}
\end{Theorem}
\marginnote{
    Integrability guarantees Lipschitz continuity and continuity guarantees differentiability.
}
Proof: \( \forall x,y \in [a,b] \), 
\begin{align}
|F(x) - F(y)| &\leq |\int_{a}^{y}f - \int_{a}^{x}f| =|\int_{x}^{y}f| \\
&\leq \int_{x}^{y} |f| \leq \left(\sup_{t \in [a,b]}|f(t)|\right) |y-x|
\end{align}
Thus \( F \) is Lipschitz continuous on \([a,b]\).

If \( f \) is continuous at \( x_0 \), then for all \( \varepsilon > 0 \), 
\(\exists \delta_\varepsilon > 0 \) such that
for all \( x \in [a,b]\), we have 
\begin{align}
    |f(x) - f(x_0)| < \varepsilon \text{ whenever } |x-x_0| < \delta_\varepsilon
\end{align}
hence for every \(t \in (-\delta_\varepsilon, \delta_\varepsilon) \setminus \{0\}\) such that \( x_0 + t \in [a,b] \), we have
\begin{align}
    \left|\frac{1}{t}\left(F(x_0 + t) - F(x_0)\right) - f(x_0)\right| &=
    \left|\frac{1}{t}\left(\int_{x_0}^{x_0+t}f\right) - \frac{1}{t}\int_{x_0}^{x_0+t}f(x_0)\right| \\
    &= \left|\frac{1}{t}\int_{x_0}^{x_0+t}(f(s) - f(x_0))ds\right| \\
    &\leq \frac{1}{|t|}\int_{x_0}^{x_0+t}|f(s) - f(x_0)|ds \\
    &\leq \frac{1}{|t|}\int_{x_0}^{x_0+t}\varepsilon = \varepsilon
\end{align}
where \(\int_{x_0}^{x_0+t} = -\int_{x_0+t}^{x_0}\) if \( t < 0 \).

It follows that, 
\begin{align}
    \limsup_{t\to 0}\left|\frac{1}{t}\left(F(x_0 + t) - F(x_0)\right) - f(x_0)\right| \leq \varepsilon
\end{align}
Letting \( \varepsilon \to 0 \) (which is allowed since 
limsup does not depends on \( t \)), we have
\begin{align}
    \lim_{t\to 0}\left|\frac{1}{t}\left(F(x_0 + t) - F(x_0)\right) - f(x_0)\right| = 0
\end{align}
\marginnote{This is how you bound things at a limit point.}
i.e, \( F \) is differentiable at \( x_0 \) and \( F'(x_0) = f(x_0) \).
\hfill \qedsymbol

\begin{Theorem}[Fundamental theorem of calculus]
    Let \( a,b \in \mathbb{R} \) be such that \( a<b \) and
    \( F: [a,b] \to \mathbb{R} \) be such that \( F \) is continuous on 
    \([a,b]\), differentiable on \( (a,b) \), and \( F' = f \in \mathcal{R}([a,b]) \).
    Then
    \begin{align}
        \int_{a}^b f(x) = F(b) - F(a)
    \end{align}
\end{Theorem}

\begin{Theorem}[Taylor's theorem with integral reminder]
    Let \( x, x_0 \in \mathbb{R} \), \( x \neq x_0 \), 
    \(I = \begin{cases}
        [x_0, x] & \text{if } x > x_0\\
        [x, x_0] & \text{if } x < x_0
    \end{cases}\), 
    \( \text{int}(I) = \begin{cases}
        (x_0, x) & \text{if } x > x_0\\
        (x, x_0) & \text{if } x < x_0
    \end{cases} \) and \( f: I \to \mathbb{R} \) be 
    n times continuously differentiable in \( I \) (i.e. 
    \( f^{(n) } \) exists and is continuous in \( I \)),
    \( n+1 \) times differentiable in \( \text{int}(I) \),
    and \( t\mapsto f^{(n+1)}(t)(x-t)^n \in \mathcal{R}(I) \).

    Then, 
    \begin{align}
        f(x) = \underbrace{\sum_{k=0}^{n}\frac{f^{(k)}(x_0)}{k!}(x-x_0)^k}_{\text{partial sum of Taylor series at \( x_0 \)}} 
        + \underbrace{\int_{x_0}^x\frac{f^{(n+1)}(t)}{n!}(x-t)^ndt}_{\text{integral reminder}}
    \end{align}
\end{Theorem}
    \marginnote{
        If \( f^{(n+1)}(t)(x-t)^n \) is not defined at the endpoints, 
        we assume it equals some Riemann integrable function at the endpoints.
    }
Proof
In the finite version of Taylor's theorem, 
we introduce \( h(t) = g(t)(x-x_0)^{n+1}-g(x_0)(x-t)^{n+1} \)
with \begin{align}
    g(t) = f(x) - \sum_{k=0}^{n}\frac{f^{(k)}(t)}{k!}(t-t)^k
\end{align}
and we proved that \( h(x) = h(x_0) = 0 \) and 
\begin{align}
    h'(t) = (n+1)g(x_0)(x-t)^n - \frac{1}{n!} f^{(n+1)}(t)(x-t)^n(x-x_0)^{n+1}
\end{align}
From the assumption on \( f \), we have that \( h \)
is continuous on \( I \), differentiable on \( \text{int}(I) \), and 
\(h'(t) \in \mathcal{R}(I)\). By applying the Fundamental 
Theorem of Calculus to \( h' \), we then obtain
\begin{align}
    0 = h(x) - h(x_0) = 
    \int_{x_0}^x h'(t) dt &= \underbrace{g(x_0)(x-x_0)^{n+1}}_{=g(x_0)\left[-(x-t)^{n+1}
    \right]_{x_0}^x} - \int_{x_0}^x\frac{f^{(n+1)}(t)}{n!}(x-t)^n(x-x_0)^{n+1}dt\\
    &= (x-x_0)^{n+1} \left(g(x_0) - \int_{x_0}^x\frac{f^{(n+1)}(t)}{n!}(x-t)^ndt\right)
\end{align}
Hence
\begin{align}
    g(x_0) = f(x) - \sum_{k=0}^{n}\frac{f^{(k)}(x_0)}{k!}(x-x_0)^k = \int_{x_0}^x\frac{f^{(n+1)}(t)}{n!}(x-t)^ndt
\end{align}
\hfill \qedsymbol

\begin{Definition}[Riemann-Stieltjes integrable]
    Let \( a \in \mathbb{R} \), \( b \in (a, \infty) \), 
    or \( b=\infty \), and \( \alpha: [a,b)\to \mathbb{R} \)
    be increasing. We say that \( f: [a,b) \to \mathbb{R} \) is
    Riemann-Stieltjes integrable with respect to \( \alpha \) on \([a,b)\) if
    \(f\in \mathcal{R}([a,s], \alpha) \) for all \( s \in (a,b) \) and the limit 
    \begin{align}
    \lim_{s\to b^-} \int_{a}^s fd\alpha
    \end{align}
    exists and if finite.
    We denote 
    \begin{align}
    \int_{a}^{b}fd\alpha \left(\text{ or } \int_{a}^{b^-}fd\alpha  \text{ if } b < \infty  \right)
    = \lim_{s\to b^-} \int_{a}^s fd\alpha
    \end{align}
    which we call R-S integral of \( f \) over \([a,b)\) with respect to \( \alpha \).

    Similarly, for \( b \in \mathbb{R} \), \( a \in (-\infty, b) \) or \( a = -\infty \), and \( \alpha: (a,b] \to \mathbb{R} \) increasing,
    we say that \( f: (a,b] \to \mathbb{R} \) is Riemann-Stieltjes integrable with respect to \( \alpha \) on \( (a,b] \) if \( f \in \mathcal{R}([s,b], \alpha) \) for all 
    \( s \in (a,b) \) and the limit
    \begin{align}
    \lim_{s\to a} \int_{s}^b fd\alpha
    \end{align}
    exists and is finite. We denote
    \begin{align}
    \int_{a}^{b}fd\alpha \left(\text{ or } \int_{a^+}^{b}fd\alpha  \text{ if } a > -\infty  \right)
    = \lim_{s\to a} \int_{s}^b fd\alpha
    \end{align}
    which we call R-S integral of \( f \) over \( (a,b] \) with respect to \( \alpha \).
\end{Definition}


\begin{Proposition}
    Let \( a, b \in \mathbb{R} \) be such that \( a<b \),
    \( \alpha: [a,b) \to \mathbb{R} \) be increasing,
    and \( f\in \mathcal{R}([a,b), \alpha) \).
    \begin{enumerate}
        \item If \( \alpha \) is continuous at \( b \) then 
        \begin{align}
            \int_{a}^{b^-}fd\alpha = \int_{a}^b fd\alpha
        \end{align}
        \item If \( \alpha \) is continuous at \( a\), then 
        \begin{align}
            \int_{a^+}^{b}fd\alpha = \int_{a}^b fd\alpha
        \end{align}
    \end{enumerate}
\end{Proposition}

\marginnote{
\begin{Remark}
If \( \alpha \) is not continuous at \( a \) or \( b \), 
then this is not true in general. For example, 
if \( \alpha =0\) on \( [a,b) \) and \( \alpha(b) = 1 \), then
\begin{align}
    \int_{a}^{b} fd\alpha = 1 \text{ but } \int_{a}^{s} fd\alpha = 0 \text{ for all } s \in (a,b)
\end{align}
\end{Remark}
}
Proof: (1.) \(\forall (a,b) \in (a,b) \), we have
\begin{align}
    \left|
        \int_{a}^{b}fd\alpha - \int_{a}^{s}fd\alpha 
    \right| =
    \left|
    \int_{s}^{b}fd\alpha
    \right| &\leq \sup_{x \in [s,b]}|f(x)| \int_{s}^{b}1d\alpha \\
    &= \sup_{x \in [s,b]}|f(x)| (\alpha(b) - \alpha(s))
\end{align}
Since \( \alpha \) is continuous at \( b \), we obtain
\( 
    \lim_{s\to b}\alpha(s) = \alpha(b)
 \), hence
\begin{align}
    \lim_{s\to b}\int_{a}^{s}fd\alpha = \int_{a}^{b}fd\alpha
\end{align}
(2.) is similar to (1.).
\hfill \qedsymbol

This allows us to define integral on functions
 that are unbounded at the endpoints, e.g. power functions.
\begin{Example}
        \begin{align}
            \int_{1}^{s}\frac{1}{x^p}dx = \begin{cases}
                \ln s - \ln 1 = \ln s \to \infty & \text{if } p=1\\
                \frac{1}{p-1}(1 - \frac{1}{s^{p-1}}) \to \begin{cases}
                \frac{1}{p-1} & \text{if } p > 1\\
                \infty & \text{if } p < 1
                \end{cases}
            \end{cases}
        \end{align}
        Hence, we have \( 
        \begin{cases}
            \int_{1}^{\infty}\frac{1}{x^p}dx = \infty & \text{if } p \leq 1\\
            x \mapsto \frac{1}{x^p} \text{ is not Riemann integrable over } [1,\infty) & \text{if } p > 1
        \end{cases}
        \).
\end{Example}

In the example we use \( \frac{d}{dx}\ln x = \frac{1}{x} \) which we proved for 
\(0<x<2\) where \( \ln(x) \) is a power series.
Proof when \(0<x<\infty\):
Observe that for all \( x,y \in (0,\infty) \), we have
\( 
    e^{\ln x + \ln y} = e^{\ln x}e^{\ln y} = xy
 \)
Hence, we have
    \( 
    \ln x + \ln y = \ln(xy)
     \)
and \(\ln x + \ln \frac1x = \ln 1 = 0\) so \( \ln \frac1x = -\ln x \).
It follows that, for every \( x \in (0,\infty) \) and \( t \) sufficiently small (\(|\frac{t}{x}| < 1\)),
\begin{align}
    \frac{1}{t}\left(
        \ln (x+t) - \ln x
    \right) &= \frac{1}{t}\ln\left(1 + \frac{t}{x}\right) \\
    &= \frac{1}{t} \sum_{k=1}^{\infty} \frac{(-1)^{k-1}}{k}\left(\frac{t}{x}\right)^k\\
    &= \sum_{k=1}^{\infty} \frac{(-1)^{k-1}}{k}\left(\frac{t}{x}\right)^{k-1} 
\end{align}
which is continuous and converges to \( \frac{1}{x} \) at \( t=0 \).
Hence, \( \ln x \) is differentiable at \( x \) and \( \frac{d}{dx}\ln x = \frac{1}{x} \).






% end of main text

\makeatletter
  \renewcommand{\section}{\@startsection{section}%
    {3}{0.8em}{-3ex \@plus -1ex \@minus -.2ex}%
    {1.5ex \@plus .2ex}
    {\hspace*{-5.5em}\fcolorbox{Periwinkle}{Periwinkle}{\parbox[c][1.0ex][b]{4em}{\phantom{space}}}
    \normalfont\Large\itshape\color{blue}}}
\makeatother

\bibliography{marginnotes}
\bibliographystyle{plainnat}

\end{document}
